\par\Needspace{9\baselineskip}\begin{center}
\rule{0.6\linewidth}{0.4pt}\par\medskip
Göttingische gelehrte Anzeigen.\quad 1814 September 26.\par\medskip
\rule{0.6\linewidth}{0.4pt}
\end{center}\pdfbookmark[1]{Göttingische gelehrte Anzeigen. 1814 September 26.}{anzeige1814}

A lecture was submitted to the Royal Society on 16 September by Professor \textsc{Gauss}, entitled:

\begin{center}
\textit{Methodus nova integralium valores per approximationem inveniendi.}
\end{center}

Among the various methods for finding approximate values of integrals, or, as the older analysts put it, for the approximate quadrature of curvilinear figures, the \textsc{Newton}-\textsc{Cotes} method, which rests on interpolation, is one of the most useful. \textsc{Newton} had given a solution to the problem of drawing a parabolic curve through any number of given points; its quadrature, which can always be carried out easily, can then serve approximately in place of the quadrature of the curve actually given through those points, and the more points are used, the more accurately it does so. \textsc{Newton}, however, had left the matter at this general indication and had merely given, as an example, for four equally spaced ordinates $A$, $B$, $C$, $D$, the approximate area between the first and last, when their distance is $=R$, as $(\frac{1}{8}A+\frac{3}{8}B+\frac{3}{8}C+\frac{1}{8}D)R$. \textsc{Cotes}, who had independently undertaken similar investigations as early as 1707, even before \textsc{Newton}'s \textit{Methodus differentialis} had appeared, was prompted by the elegant form in which \textsc{Newton} had presented the final result in the above example (\textit{pulcherrima et}
\SourceLeaf{212}
\textit{utilissima regula}, as \textsc{Cotes} calls it) to extend these rules further, up to the case of 11 ordinates. The required area thus always appears as the product of the base, or distance between the outermost ordinates, and the sum of the ordinates multiplied by definite numerical coefficients; two ordinates equally distant from the beginning and the end always have equal coefficients. \textsc{Cotes} gives these quadrature coefficients up to the case of 11 ordinates at the end of the paper \textit{de methodo differentiali}, which forms part of the \textit{Harmonia mensurarum}, without further explanation of how he computed them. Perhaps this unassuming brevity, which presents only the final result, accounts for this beautiful and useful method seeming to be less known and used by analysts than it deserves.

This method rests throughout on the assumption of equal distances between the ordinates. At first sight this assumption certainly seems the simplest and most natural, and the question had not yet been raised whether, despite this, it might be more advantageous to take ordinates at unequal distances. To decide this question, the theory of the quadrature coefficients first had to be developed in unrestricted generality, and the degree of accuracy of the result determined. It appeared that the conditions on which this degree of accuracy depends are such that it can indeed be doubled by suitably choosing ordinates at unequal distances, so that any number of appropriately chosen ordinates goes as far as twice that number at equal distances. These investigations, together with the complete theory of the most suitable choice of ordinates, the quadrature coefficients to be used with them and the determination of the degree of accuracy afforded by this procedure, form the main contents of the present paper.

From the brief development of the theory of the \textsc{Cotes} quadrature coefficients, which the author thought it necessary to put first, we touch here only on what concerns the degree of accuracy of the approximate integrals found by it. Above all, it must be noted that the applicability of this method, like that of interpolation, rests on the assumption that the ordinates within the region whose quadrature is required
\SourceLeaf{213}
can be represented by a convergent series. Let \(x\) be the abscissa, \(y\) the ordinate, and let the integral \(\int y\,\mathrm{d}x\) be required from \(x=g\) to \(x=h\). Introduce another variable in place of \(x\) by setting, for example, \(x=g+(h-g)t\), or alternatively \(x=\frac12(g+h)+\frac12(h-g)u\). Here, then, \(y\) must be representable by series such as
\[
\alpha+\alpha't+\alpha''tt+\alpha'''t^3+\text{etc.}
\]
or
\[
\beta+\beta'u+\beta''uu+\beta'''u^3+\text{etc.}
\]
which converge as long as \(t\) in the former case, at least, and \(u\) in the latter do not become greater than 1. For brevity, the coefficients \(\alpha'\) and \(\beta'\) may therefore be assigned order 1, the coefficients \(\alpha''\) and \(\beta''\) order 2, and so on. This being assumed, it is shown that the errors incurred in the \textsc{Cotes} method are always of higher order as the number of values of \(y\) used increases, but in such a way that an odd number and the next larger even number always produce errors of the same order. Thus, for three ordinates the error is very nearly \(=\frac{1}{120}(h-g)\alpha''''\), for four ordinates nearly \(=\frac{1}{270}(h-g)\alpha''''\); then, for five ordinates nearly \(=\frac{1}{2688}(h-g)\alpha^{\mathrm{VI}}\) and for six ordinates nearly \(=\frac{11}{52500}(h-g)\alpha^{\mathrm{VI}}\), and so on. From this one sees that, in general, it will be advantageous to use an odd number of ordinates when applying the \textsc{Cotes} method.

The author then turns to the general investigation in which the restriction that the ordinates be equally spaced is removed. If \(A,A',A''\), etc., are the values of \(y\), while the corresponding values of \(t\) are \(a,a',a''\), etc., or \(b,b',b''\), etc., are the corresponding values of \(u\), and their number is \(n+1\), the approximate integral will again have the form
\[
(h-g)(RA+R'A'+R''A''+\text{etc.})
\]
where \(R,R',R''\), etc., are numerical coefficients which, independently of the function \(y\), are determined solely by \(a,a',a''\), etc., or by \(b,b',b''\), etc. The author's investigations give the following result for determining them. Let
\[
T=(t-a)(t-a')(t-a'')\,\ldotp\,\ldotp\,\ldotp\,\ldotp
\]
\SourceLeaf{214}
Multiplication of this polynomial in \(t\), which rises to degree \(n+1\), by the infinite series representing the logarithm of \(\frac{t}{t-1}\), namely
\[
t^{-1}+\frac12t^{-2}+\frac13t^{-3}+\text{etc.}
\]
is to give the product \(T'+T''\), so that \(T'\) denotes the polynomial in \(t\) contained in it, and \(T''\) the remaining series extending to infinity in negative powers of \(t\). This being assumed, the quadrature coefficients \(R,R',R''\), etc., are obtained by substituting in \(\frac{T'\mathrm{d}t}{\mathrm{d}T}\), for \(t\), successively the values \(a,a',a''\), etc. In a similar and somewhat more convenient way these coefficients are derived from \(b,b',b''\), etc., by denoting the function
\[
(u-b)(u-b')(u-b'')\,\ldotp\,\ldotp\,\ldotp\,\ldotp
\]
by \(U\), its product with the infinite series
\[
u^{-1}+\frac13u^{-3}+\frac15u^{-5}+\text{etc.}
\]
by \(U'+U''\) (so that \(U'\) represents the polynomial in \(u\) contained in it), and then substituting in \(\frac{U'\mathrm{d}u}{\mathrm{d}U}\), for \(u\), successively the values \(b,b',b''\), etc. In place of the rational functions \(\frac{T'\mathrm{d}t}{\mathrm{d}T},\frac{U'\mathrm{d}u}{\mathrm{d}U}\), polynomials in \(t\) and \(u\) can also be found which can replace them, and for whose determination the author develops a general method.

The degree of accuracy of the integration formula now depends on the nature of the series \(T''\) or \(U''\). In general the error is indeed of order \(n+1\); but if some of the first terms of these series vanish, the error is of higher order, so that, if \(T''\) begins only with the power \(t^{-m}\), or \(U''\) with the power \(u^{-m}\), the error is of order \(n+m\).

It followed that, insofar as the values \(a,a',a''\), etc., or \(b,b',b''\), etc., can be chosen arbitrarily, it must be possible to determine them so that the first \(n+1\) terms of \(T''\) or \(U''\) actually vanish; the consequence will be that the error of the integration formula is of order \(2n+2\). The investigation therefore proceeds to determine those functions \(T\) and \(U\), for each value of \(n\), which satisfy the stated condition. Limited space does not permit
\SourceLeaf{215}
us to enter into the details of this investigation here: we therefore remark only that these functions follow a very simple law of progression and are closely connected with the expansion of the series
\[
\begin{aligned}
&t^{-1}+\tfrac{1}{2}t^{-2}+\tfrac{1}{3}t^{-3}+\tfrac{1}{4}t^{-4}+\mathit{etc.}\\
&u^{-1}+\tfrac{1}{3}u^{-3}+\tfrac{1}{5}u^{-5}+\tfrac{1}{7}u^{-7}+\mathit{etc.}
\end{aligned}
\]
into continued fractions, which the author gave in an earlier paper [\textit{Disquisitiones generales circa seriem infinitam}
\(1+\frac{\alpha.\beta}{1.\gamma}x+\mathit{etc.}\)].\,---
Solving the equation \(T=0\) or \(U=0\) then evidently gives the values of \(a,a',a''\), etc., or \(b,b',b''\), etc., and the quadrature coefficients are determined by the general rules, which admit further special simplifications in this case. In most cases, however, both the values of \(a,a',a''\), etc., and the quadrature coefficients will indeed be irrational quantities. This is in itself quite immaterial, provided only that their numerical values have been computed once and for all to an appropriate degree of accuracy. Once this is done, applying this method to any number of ordinates will require little or no more effort than applying the \textsc{Cotes} method to the same number; the latter, however, would have to be applied to twice as many to give approximately the same accuracy of the result as the former.

\Needspace{17\baselineskip}To leave nothing to be desired for the application of this new method, the author has also supplied the numerical values of \(a,a',a''\), etc., and of \(R,R',R''\), etc., computed to 16 decimal places, together with the \textsc{Brigg}sian logarithms of the latter to 10 decimal places, all up to the case of seven ordinates. In this last case the error of the integration formula will be nearly \(=\frac{1}{176679360}\alpha^{\mathrm{XIV}}\), from which one may judge that, in most cases occurring in practice, a smaller number will suffice. To make the application of the rules and their relative accuracy still more evident, the computation of \(\int\frac{\mathrm{d}x}{\log x}\) from \(x=100000\) to \(x=200000\) is appended as an example. Even with four values, the error is only \(\frac{1}{56000000}\) of the whole; with a larger number it is lost among the unavoidable errors still left even by the larger logarithm tables.
\par\begin{center}\rule{0.22\linewidth}{0.3pt}\end{center}
